\clearpage
\section{Vorticity}
\begin{colbox}
	\begin{enumerate}
		\item Write the definition of \(\omega_{ij}\) and multiply both sides by \(\epsilon_{ijl}\). Simplify to obtain a single summand.
		\item Multiply the original definition by \(\epsilon_{ijl}\) and use Levi-Civita tensor identities to obtain \begin{equation}
			\vb{\omega} = \curl{\vb{v}}
		\end{equation}
	\end{enumerate}
\end{colbox}
\subsection{part 1}
The definition of the vorticity is 
\begin{equation}
	\omega_{ij} = -\frac{1}{2}\epsilon_{ijk}\omega_k
\end{equation}
multiplying by \(\epsilon_{ijl}\)
\begin{align}
	\epsilon_{ijl}\omega_{ij} &= -\frac{1}{2}\epsilon_{ijl}\epsilon_{ijk}\omega_k\\
	\epsilon_{ijl}\omega_{ij} &= -\frac{1}{2}\ins{2\delta_{kl}\cdot\ins{-1}^{2}}\omega_k\\
	\epsilon_{ijl}\omega_{ij} &= -\omega_l
\end{align}
\subsection{part 2}
By original definition I imagine this is what's meant:
\begin{equation}
	\omega_{ij} = \frac{1}{2}\ins{\pdv{v_i}{x_j}-\pdv{v_j}{x_i}}
\end{equation}
multiplying both sides by \(\epsilon_{ijl}\)
\begin{align}
	\epsilon_{ijl}\omega_{ij} &= \frac{1}{2}\epsilon_{ijl}\ins{\pdv{v_i}{x_j}-\pdv{v_j}{x_i}} \\
	&= \frac{1}{2}\ins{\epsilon_{ijl}\pdv{v_i}{x_j}-\epsilon_{ijl}\pdv{v_j}{x_i}}\\
	&= -\epsilon_{ijl}\pdv{v_i}{x_j}
\end{align}
combining with the result of part 1 we get
\begin{equation}
	\omega_l = \epsilon_{ijl}\pdv{v_i}{x_j}
\end{equation}
which by definition is 
\begin{equation}
	\vb{\omega} = \curl{\vb{v}}
\end{equation}